\section{Numerical results}
\label{sec:numerical}
This section presents numerical experiments designed to validate the theoretical analysis and evaluate the practical performance of the proposed algorithm. 
Specifically, we verify its core properties including linear convergence, bound-preserving property, and energy stability in \Cref{subsec:test}. 
Numerical simulations in two-dimensional and three-dimensional settings are demonstrated in \Cref{subsec:2d} and \Cref{subsec:3d}, respectively, 
illustrating the effectiveness and robustness of our solver.


\subsection{Convergence test}
\label{subsec:test}
In this numerical experiment, the parameters are set as follows:
\begin{equation}
    L=2\pi, \epsilon=0.10, \tau=10^{-4}, T=0.1, \theta=4.0.
\end{equation}
The initial condition is
\begin{equation}
    u_0(x,y)=0.5+0.25\sin(x)\sin(y).
\end{equation}
\par For the iterative solver, the penalty parameter $\rho$ is dynamically updated based on the primal and dual residuals to ensure balanced convergence. \\
The convergence criterion is set as
\begin{equation}
    \max\{r,s\}\leq \gamma=10^{-8},
\end{equation}
where $r=\|u_1^{(k)}-u_2^{(k)}\|_h$ and $s=\|u_1^{(k)}-u_1^{(k-1)}\|_h$ denote the primal residual and the dual residual, respectively.
\par The numerical solution obtained with $N=512,\tau=10^{-4}$ is taken as the exact solution, and the spatial convergence rate is tested for $N=16,32,64,128,256$.
The temporal convergence order test is analogous.
The test results, displayed in \Cref{ta:rate}, confirm the second order in space and first order in time for our solver.
\par Fixing $N=256$ and $\tau = 0.01$, we show the linear convergence and other properties of the algorithm in \Cref{fig:1}.
From \Cref{fig:11}, we see the number of ADMM iterations at each time step, which speaks to the algorithm's effectiveness. 
The discrete energy, plotted in \Cref{fig:12}, decays monotonically as expected, confirming energy stability.
\Cref{fig:13} and \Cref{fig:14} traces the maximum and minimum of $u$ over time, both staying safely within the physical bounds $(0,1)$, demonstrating the bound-preserving property. 
\Cref{fig:15} displays the primal residual during the second time step; 
its linear decay on a semi-log scale indicates a linear convergence rate.
Finally, \Cref{fig:16} captures the dynamial evolution process of the phase variable at different time steps.
\begin{table}
    \centering
    \begin{tabular}{|c|c|c||c|c|c|}
    \hline
    \multicolumn{3}{|c||}{Spatial convergence}  & \multicolumn{3}{|c|}{Temporal convergence} \\
    \hline
    $h$&Error&Rate&$\tau$&Error&Rate\\
    \hline
    2$\pi$/16&3.026E-2&-&0.02&4.336E-2&-\\
    \hline
    2$\pi$/32&7.635E-3&1.99&0.02/2&2.246E-2&0.95\\
    \hline
    2$\pi$/64&1.918E-3&1.99&0.02/4&1.086E-2&1.05\\  
    \hline
    2$\pi$/128&4.853E-4&1.98&0.02/8&4.741E-3&1.20\\
    \hline
    2$\pi$/256&1.282E-4&1.92&0.02/16&1.595E-3&1.57\\
    \hline

    \end{tabular}
    \caption{Errors and convergence rates.}
    \label{ta:rate}
\end{table}


\begin{figure}[htbp]
    \centering
    \begin{subfigure}[b]{0.49\textwidth}
        \centering
        \includegraphics[width=\textwidth]{figure/admm_iterations_1.pdf}
        \caption{The number of iterations over time $t$.}
        \label{fig:11}
    \end{subfigure}
    \hfill
    \begin{subfigure}[b]{0.49\textwidth}
        \centering
        \includegraphics[width=\textwidth]{figure/energy_curve_1.pdf}
        \caption{The energy over time $t$.}
        \label{fig:12}
    \end{subfigure}
    
    \vspace{0.5cm}
    
    \begin{subfigure}[b]{0.49\textwidth}
        \centering
        \includegraphics[width=\textwidth]{figure/bounds_upper_1.pdf}
        \caption{$1-\max (u)$ over time $t$.}
        \label{fig:13}
    \end{subfigure}
    \hfill
    \begin{subfigure}[b]{0.49\textwidth}
        \centering
        \includegraphics[width=\textwidth]{figure/bounds_lower_1.pdf}
        \caption{$\min (u)$ over time $t$}
        \label{fig:14}
    \end{subfigure}
    \vspace{0.5cm}

    \begin{subfigure}[b]{0.49\textwidth}
        \centering
        \includegraphics[width=\textwidth]{figure/admm_convergence_1.pdf}
        \caption{The primal residual in the second time step.}
        \label{fig:15}
    \end{subfigure}
    
    \caption{Convergence test}
    \label{fig:1}
\end{figure}

\begin{figure}[htbp]
    \centering
    \includegraphics[width=\textwidth]{figure/admm_evolution_1.pdf}
    \caption{The dynamial evolution over time $t$.}
    \label{fig:16}
    \hfill
\end{figure}


\subsection{Two-dimensional simulation}
\label{subsec:2d}
In this numerical experiment, we show a two-dimensional simulation process here.
\par The following parameters are adopted in this example:
\begin{equation}
L=2.0, N=512,\epsilon=0.05, \tau=0.1, T=100.
\end{equation}
The parameter $\theta$ is chosen as the variable parameter in this experiment, which is taken as $\theta=3.0,4.0,5.0$. 
Notice that the time step $\tau=0.1$ is significantly large, 
which verifies that our method has no requirement for the time step.
\par The initial condition is
\begin{equation}
    u_0(x,y)=0.01+0.98\text{rand}(x,y).
\end{equation}
where $\text{rand}(x, y)$ is the function producing the random numbers in $(0,1)$. 
\par We set the convergence criterion for \Cref{alg:ADMM} as:
\begin{equation}
\max\{r,s\}\leq \gamma=10^{-8}.
\end{equation}

For $\theta=3.0$, the number of iterations, the energy, $1-\max (u)$ and $\min (u)$ over time $t$ are given in \Cref{fig:2_1}.
The dynamic evolution of the phase variable is depicted in \Cref{fig:2_2} for different choices of $\theta$.


\begin{figure}[htbp]
    \centering
    \begin{subfigure}[b]{0.49\textwidth}
        \centering
        \includegraphics[width=\textwidth]{figure/admm_iterations_2.pdf}
        \caption{The number of iterations over time $t$.}
        \label{fig:21}
    \end{subfigure}
    \hfill
    \begin{subfigure}[b]{0.49\textwidth}
        \centering
        \includegraphics[width=\textwidth]{figure/energy_curve_2.pdf}
        \caption{The energy over time $t$.}
        \label{fig:22}
    \end{subfigure}
    
    \vspace{0.5cm}
    
    \begin{subfigure}[b]{0.49\textwidth}
        \centering
        \includegraphics[width=\textwidth]{figure/bounds_upper_2.pdf}
        \caption{$1-\max (u)$ over time $t$.}
        \label{fig:23}
    \end{subfigure}
    \hfill
    \begin{subfigure}[b]{0.49\textwidth}
        \centering
        \includegraphics[width=\textwidth]{figure/bounds_lower_2.pdf}
        \caption{$\min (u)$ over time $t$}
        \label{fig:24}
    \end{subfigure}
    \vspace{0.5cm}
    
    \caption{2D simulation}
    \label{fig:2_1}
\end{figure}



\begin{figure}[htbp]
    \centering
    \begin{subfigure}[b]{0.78\textwidth}
        \centering
        \includegraphics[width=\textwidth]{figure/admm_evolution_2_3.0.pdf}
        \caption{$\theta=3.0$}
        \label{fig:25}
    \end{subfigure}
    \hfill
    
    \vspace{0.5cm}
    
    \begin{subfigure}[b]{0.78\textwidth}
        \centering
        \includegraphics[width=\textwidth]{figure/admm_evolution_2_4.0.pdf}
        \caption{$\theta=4.0$}
        \label{fig:26}
    \end{subfigure}
    \hfill

    \vspace{0.5cm}

    \begin{subfigure}[b]{0.78\textwidth}
        \centering
        \includegraphics[width=\textwidth]{figure/admm_evolution_2_5.0.pdf}
        \caption{$\theta=5.0$}
        \label{fig:27}
    \end{subfigure}
    
    \caption{2D simulation: the dynamial evolution over time $t$.}
    \label{fig:2_2}
\end{figure}

\iffalse 
\subsection{long simulation}\label{subsec:Adaptive}
In this experiment, we adopt an adaptive time-stepping strategy to demonstrate the high efficiency of the algorithm. 
The parameters we adopted are set as
\begin{equation}
L=2.0, N=256,\epsilon = 0.05, \tau^0=1.0,\theta=4.0, T=10^{12}.
\end{equation}
\par The initial condition is
\begin{align}
u_0(x,y)=
\begin{cases}
    10^{-5},&\text{if}\quad |x-1|\leq 0.35, |y-1|\leq 0.35,\\
    1-10^{-5},&\text{otherwise}
\end{cases}
\end{align}
which is taken from \cite{wang2020stabilized}.
Numerical experiments indicate that the two-phase evolution under this initial condition requires long-time simulation. 
Hence, we employ this example to perform adaptive time-stepping simulation, 
which dynamically adjusts the time step based on the energy decrement to maintain accuracy while reducing computational cost.
To be more precise, the adaptive adjustment procedure is:
\begin{align}
\tau^{n+1}=
\begin{cases}
\min \{\beta_1 \tau^{n},\tau_{\max}\},&\text{if} \quad |\frac{E_h^{n-1}-E_h^n}{E_h^{n-1}}|<\gamma_1,\\
\max \{\beta_2 \tau^{n},\tau_{\min}\},&\text{if} \quad |\frac{E_h^{n-1}-E_h^n}{E_h^{n-1}}|>\gamma_2,\\
\tau^n,&\text{otherwise},
\end{cases}
\end{align}
where $\beta_1=2.0,\tau_{\max}=10^{10},\gamma_1=0.1,\beta_2=0.5,\tau_{\min}=10^{-3},\gamma_2=0.2$.
\par The numerical results are shown in \Cref{fig:4}. 
The same observations as in \Cref{subsec:test} hold here, that is, 
the results confirm linear convergence, energy stability, and the bound-preserving property of the scheme from \Cref{fig:41} to \Cref{fig:44}.
The evolution of the phase field variable over time is shown in \Cref{fig:45}, where the square-shaped phase is seen to rapidly shrink into a circular one.

\begin{figure}[htbp]
    \centering
    \begin{subfigure}[b]{0.49\textwidth}
        \centering
        \includegraphics[width=\textwidth]{figure/admm_iterations_adaptive.png}
        \caption{The number of iterations}
        \label{fig:41}
    \end{subfigure}
    \hfill
    \begin{subfigure}[b]{0.49\textwidth}
        \centering
        \includegraphics[width=\textwidth]{figure/bounds_preservation_adaptive.png}
        \caption{Bound preservation}
        \label{fig:42}
    \end{subfigure}
    
    \vspace{0.5cm}
    
    \begin{subfigure}[b]{0.49\textwidth}
        \centering
        \includegraphics[width=\textwidth]{figure/energy_curve_adaptive.png}
        \caption{Energy stability}
        \label{fig:43}
    \end{subfigure}
    \hfill
    \begin{subfigure}[b]{0.49\textwidth}
        \centering
        \includegraphics[width=\textwidth]{figure/admm_convergence_adaptive.png}
        \caption{Primal residual}
        \label{fig:44}
    \end{subfigure}
    \vspace{0.5cm}

    \begin{subfigure}[b]{1.00\textwidth}
        \centering
        \includegraphics[width=\textwidth]{figure/admm_evolution_test_4.png}
        \caption{Dynamical evolution}
        \label{fig:45}
    \end{subfigure}
    
    \caption{Adaptive time-stepping simulation}
    \label{fig:4}
\end{figure}

\fi


\subsection{Three-dimensional simulation}
\label{subsec:3d}
We present a numerical simulation experiment in a three-dimensional setting here. The parameters we adopted are 
\begin{equation}
L=1.0, N=64,\theta=4.0, \tau=0.01, T=0.5.
\end{equation}
The parameter $\epsilon$ is chosen as the variable parameter in this experiment, which is taken as $\epsilon=0.05,0.10,0.15$.
\par The initial condition is
\begin{equation}
    u_0(x,y,z)=0.45+0.10\text{rand}(x,y,z).
\end{equation}
where $\text{rand}(x,y,z)$ is the function producing the random numbers in $(0,1)$. 
\par The convergence criterion is:
\begin{equation}
\max\{r,s\}\leq \gamma=10^{-8}.
\end{equation}

\begin{figure}[htbp]
    \centering
    \begin{subfigure}[b]{0.49\textwidth}
        \centering
        \includegraphics[width=\textwidth]{figure/admm_iterations_3.pdf}
        \caption{The number of iterations over time $t$.}
        \label{fig:31}
    \end{subfigure}
    \hfill
    \begin{subfigure}[b]{0.49\textwidth}
        \centering
        \includegraphics[width=\textwidth]{figure/energy_curve_3.pdf}
        \caption{The energy over time $t$.}
        \label{fig:32}
    \end{subfigure}
    
    \vspace{0.5cm}
    
    \begin{subfigure}[b]{0.49\textwidth}
        \centering
        \includegraphics[width=\textwidth]{figure/bounds_upper_3.pdf}
        \caption{$1-\max (u)$ over time $t$.}
        \label{fig:33}
    \end{subfigure}
    \hfill
    \begin{subfigure}[b]{0.49\textwidth}
        \centering
        \includegraphics[width=\textwidth]{figure/bounds_lower_3.pdf}
        \caption{$\min (u)$ over time $t$}
        \label{fig:34}
    \end{subfigure}
    \vspace{0.5cm}
    
    \caption{3D simulation}
    \label{fig:3_1}
\end{figure}


\begin{figure}[htbp]
    \centering
    \begin{subfigure}[b]{0.78\textwidth}
        \centering
        \includegraphics[width=\textwidth]{figure/admm_3d_evolution_3_0.05.pdf}
        \caption{$\epsilon=0.05$}
        \label{fig:35}
    \end{subfigure}
    \hfill
    
    \vspace{0.5cm}
    
    \begin{subfigure}[b]{0.78\textwidth}
        \centering
        \includegraphics[width=\textwidth]{figure/admm_3d_evolution_3_0.10.pdf}
        \caption{$\epsilon=0.10$}
        \label{fig:36}
    \end{subfigure}
    \hfill

    \vspace{0.5cm}

    \begin{subfigure}[b]{0.78\textwidth}
        \centering
        \includegraphics[width=\textwidth]{figure/admm_3d_evolution_3_0.15.pdf}
        \caption{$\epsilon=0.15$}
        \label{fig:37}
    \end{subfigure}
    
    \caption{3D simulation: the dynamial evolution over time $t$.}
    \label{fig:3_2}
\end{figure}
For $\epsilon=0.05$, the number of iterations, the energy, $1-\max (u)$ and $\min (u)$ over time $t$ are given in \Cref{fig:3_1}. 
\Cref{fig:3_2} shows the interface of $u=0.50$ by interpolating the numerical solution at different time steps from the three-dimensional simulation. 
The results indicate that the algorithm accurately captures the dynamical process of the two-phase interface with satisfactory precision,
which is consistent with the physical phenomenon of phase separation. In summary, the numerical result demonstrates the robustness of our solver in numerical simulation.



